% Part: propositional-logic
% Chapter: propositional-logic 
% Section: soundness

\documentclass[../../../include/open-logic-section]{subfiles}

\begin{document}

\olfileid{pl}{prp}{snd}

\olsection{د قضيوي منطق سموالے}

\begin{thm}[سموالے]
\ollabel{thm:soundness}
که $\Gamma\Proves!A$ نو $\Gamma \Entails !A$۔
\end{thm}

\begin{proof}
د اشتقاق پر جوړښت استقرا کوو۔ که $!A$ بديهي اصل وي، نو $\Entails
!A$ او له دې څخه $\Gamma \Entails!A$ هم راځي۔ همداسې که $!A \in \Gamma$، نو $\Gamma \Entails!A$۔ په استقرايي ګام کښې فرض کړئ $!A$ له $!C$ او $!C\lif!A$ څخه په \emph{مودوس پوننس} اخيستل شوے دے۔ د استقرا د فرضيې له مخې قضيه د $!C\lif !A$ او $!C$ دپاره سمه ده۔ د \olref[pl][syn][sem]{prop:semanticalfacts} دويم توکے مطلوبه پايله ورکوي۔
\end{proof}

\begin{cor}
که $\Gamma$ د صدق وړ وي، نو $\Gamma$ سازګار دے۔ په نتيجه کښې قضيوي منطق سازګار دے۔
\end{cor}

\end{document}
